% ============================================================================
% Section 5 -- implementation.
%
% COMPRESSED 2026-09-27. The tier-aware-attention and scheduler subsections
% restated section 4's fig:exec and admission rules almost verbatim; in a 10pp
% body that duplication is the first thing to go. What is left is what only the
% implementation can tell you: the two failure modes a design sketch omits, the
% handover cost, and the composition of F -- the number the result turns on.
%
% PROVENANCE: these are engine timings taken during development, reproducible
% from the harnesses in paper/GH200_RUNBOOK.md but not archived as JSON. Said
% once here rather than as a "% src:" per number pointing at a commit.
% ============================================================================

\section{Implementation}
\label{sec:impl}

The system is implemented in vLLM~v1 in roughly 3{,}900 lines across 22 files.
Every path is disabled by default and selected at runtime, so a build with the
feature off exercises the stock code, which \S\ref{sec:eval:setup} verifies.
Timings here are engine measurements taken during development on the same A100-80,
reproducible from the harnesses in our artifact but not part of the frozen dataset
behind \S\ref{sec:eval}.

\paragraph{Host pool and attention.}
\label{sec:impl:pool}
The pool owns a private block-id space, a per-request map from logical to host
block, and two pinned host tensors per layer (\S\ref{sec:design:spaces}).
Attention subclasses the FlashAttention~\cite{flashattention2} backend: host-tier
rows are gathered with \texttt{index\_select}, attended against the rolling slot
with the compacted block table, and scattered back with \texttt{index\_copy\_}.

\paragraph{Two failure modes a design sketch omits.} Both were silent rather than
fatal, which is why they are worth reporting. First, deferring a block-starved
request at the head of the waiting queue is \emph{separately} necessary: admission
with an empty block list and a prompt marked already computed are jointly
insufficient, because a request needing no blocks can still be starved behind one
that does. Second, block accounting is length-based, so a nulled entry still counts
as present --- \texttt{allocate\_slots} never refills it, and a promotion gate asking
``do the blocks cover the whole sequence?'' passes on an all-null table and serves
garbage. Refill is therefore all-or-nothing, the gate counts non-null entries, and
refilled ids reach the worker through a rewrite channel because they replace
positions rather than append.

\paragraph{Prefill handover.}
\label{sec:impl:stash}
A host-tier request's KV must reach the host pool before admission. We implement a
prompt-keyed stash whose entry outlives the producing request, with admission
contingent on the worker acknowledging the KV is present. It costs \textbf{0.22\,ms
per decode step} at 2048 context (14.2\,ms over 64 steps, 38\,GB/s), about 3\% of a
host-tier request's per-step cost.

\paragraph{What the host path costs per step.}
\label{sec:impl:cost}
Beyond the transfer, using the host path costs a fixed \textbf{3.4\,ms} per step at
64 layers --- the $F$ of \S\ref{sec:model:omits}, on which the paper's result turns.
We previously attributed it to the per-layer \emph{issue structure}: an operator
trace showed no dominant term, just some 640 extra kernel launches per step, ten per
layer at $0.19$\,ms a layer. That attribution was wrong, and testing it is what
found the real cost. A measurement-only mode that issues whole extents, cutting the
launch count eightfold, moved neither $F$ nor $c$ by more than their run-to-run
drift. $F$ is not launch-bound.

Ablating the mechanisms instead puts $1.21$\,ms of it in the extra attention call
that the two-source design requires (\S\ref{sec:design:exec}) and the other
$2.15$ in the mirror write and the remaining per-layer work. $F$ scales with layer
depth at $0.053$\,ms per layer, and it is \emph{not} machine-independent, so a Grace
host may pay differently --- which is why measuring it on a coherent-link part is
the first experiment we ask of one. Appendix~\ref{app:ridercost} gives the full
ledger of what was tried and what landed.
